\scanpage{31}
\mspage{12}
qui implique bien sûr (26).

\margin{\unsure{revenant} \unsure{avant} l'introduction de $u$}
Je \deleted{vais} considère les isomorphismes
\diagram{\begin{tikzcd}[ampersand replacement=\&,column sep=large,row sep=large]
\Pi_{0,3}\arrow[r,"\widetilde\psi_i^\omega"]\arrow[dr,"\psi_i^\omega"']\&\pi(P^\omega)\\
\&\pi(Q_i)\arrow[u,"(d_{\unsure{i}}^\omega)^{-1}"']
\end{tikzcd}}\quad(38)
\margin{\uline{NB} le passage de $u$ à $\tau(u)$ s'exprime par
$\xi\mapsto\sigma_\infty\tau_\infty(\xi)=\widetilde\sigma_\infty(\xi)$}
\margin{\deleted{\unsure{Changement} \unsure{de base} de l'iso $\Pi_{0,3}\simeq\Pi_{0,3}$ pour \unsure{compatib.}}}

de sorte que les relations (35) et (19) donnent
\[
(39)\qquad \widetilde\psi_{\omega i}^\omega(g)=\widetilde\psi_i^\omega(\rho g)
\qquad\text{i.e.}\quad\widetilde\psi_{\omega i}=\widetilde\psi_i\circ\rho,
\]
\inserted{\figure{\begin{tikzpicture}\node[draw,ellipse] {$\rho^{-1}$?};\end{tikzpicture}}}
et j'introduis les chemins
\[
(40)\qquad\widetilde d_i^\omega:
\begin{tikzcd}[ampersand replacement=\&,column sep=large]
P_{\unsure{\infty|\xi}}^\omega\arrow[r,no head]\&P_{\unsure{\infty|\xi}}^{\omega'}
\end{tikzcd},\qquad
\widetilde d_i^\omega=(d_i^{\omega'})^{-1}d_i^\omega
\]
permutés de façon simplement transitive par $G_J$,
\deleted{\unsure{de} \unsure{\illegible{1}} et satisfaisant} satisfaisant
\[
(41)\qquad\tau(\widetilde d_i^\omega)=\widetilde d_i^{\omega'},\qquad
 d_i^\omega d_i^{\omega'}=1
\]
On \deleted{\unsure{aura donc}} \inserted{a} les relations
\[
(42)\qquad\widetilde d_i^\omega(\widetilde\psi_i^\omega(g))=
\widetilde\psi_i^{\omega'}(\sigma_\infty(g)).
\]
Revenant à l'\deleted{auto}morphisme $u$, on aura \deleted{automatiquement}
\[
(43)\qquad u(\widetilde d_i^\omega)=\widetilde\psi_i^{\omega'}(\unsure{\sigma_\infty}(\widetilde\beta))\widetilde d_i^\omega
=\widetilde d_i^\omega\widetilde\psi_i^\omega(\widetilde\beta)
\]
avec
\[
(44)\qquad\widetilde\beta=\xi\sigma_\infty(\xi)^{-1}.
\]
satisfaisant \inserted{donc} \deleted{\unsure{ainsi}} comme $\widetilde\beta$ doit, la relation
\[
(45)\qquad\widetilde\beta\sigma_\infty(\widetilde\beta)=1.
\]

Considérons maintenant \unsure{le|les} ss-\unsure{groupoïde|groupes} de
\margin{30}
